%====================================
% KAPITEL 1: EINLEITUNG
%====================================
\section{Einleitung}

\subsection{Problemstellung und Relevanz}

Bereits im Jahr 2003 ver\"offentlichten DeLone und McLean mit ihrem aktualisierten Information-Systems-Success-Modell einen der bis heute einflussreichsten Ans\"atze zur Bewertung des Erfolgs von Informationssystemen. Das Modell beschreibt den Erfolg technologischer Systeme anhand mehrerer miteinander verkn\"upfter Dimensionen und gilt als zentraler Referenzrahmen der Informationssystemforschung.\footcite[Vgl.][S.~9--10]{delone2003} Trotz dieses etablierten Verst\"andnisses davon, was den Erfolg eines Systems ausmacht, bleiben ERP-Implementierungen in der Praxis jedoch h\"aufig hinter den Erwartungen zur\"uck. Aktuelle systematische Literaturanalysen zeigen, dass zahlreiche Unternehmen die angestrebten Potenziale ihrer ERP-Systeme selbst nach der Einf\"uhrung nicht vollst\"andig realisieren.\footcite[Vgl.][S.~2--3]{butarbutar2023}

Dieser Befund verdeutlicht, dass technologische Exzellenz allein nicht ausreicht. Schon zwei Jahre zuvor hatten Nah, Lau und Kuang die elf kritischen Erfolgsfaktoren f\"ur ERP-Projekte systematisiert und dabei das Change Management explizit als eigenst\"andigen Faktor identifiziert.\footcite[Vgl.][S.~285]{nah2001} Die Forschung zeigt, dass Change Management und die Unterst\"utzung durch das Top-Management zu den zentralen Erfolgsfaktoren von ERP-Implementierungen z\"ahlen.\footcites[Vgl.][S.~1181--1190]{leyh2014}[][S.~258--259]{somers2004} Diese Erfolgsproblematik ist nicht neu: Schon Holland und Light wiesen darauf hin, dass rein technologische Ans\"atze ohne Ber\"ucksichtigung der organisatorischen Passung oft defizit\"ar bleiben.\footcite[Vgl.][S.~30--36]{holland1999}

Um diesen Herausforderungen zu begegnen, bietet das Change Management die notwendige theoretische sowie praktische Perspektive. W\"ahrend Armenakis und Bedeian die Grundlagen f\"ur den Umgang mit organisationalem Wandel schufen,\footcite[Vgl.][S.~300--302]{armenakis1999} verdeutlicht die Forschung von Hubbart, dass insbesondere die individuelle Abwehrhaltung gegen\"uber Ver\"anderungen eine erhebliche H\"urde bleibt.\footcite[Vgl.][S.~162]{hubbart2023} Den \"Ubergang zwischen allgemeinen Strategien und dem spezifischen ERP-Kontext bildet der Ansatz von Aladwani, der betont, dass der Projekterfolg wesentlich von einem strukturierten, prozessorientierten Umgang mit Nutzerwiderstand abh\"angt.\footcite[Vgl.][S.~266--275]{aladwani2001} Als zentrale Wirkmechanismen gelten dabei die Hebel Kommunikation, F\"uhrung und Schulung.\footcites[Vgl.][S.~285]{nah2001}[][S.~7--8]{butarbutar2023} Obwohl die Relevanz dieser Faktoren in der Forschung breit belegt ist, fokussiert der Gro\ss{}teil der Studien bislang Gro\ss{}unternehmen und Konzerne.\footcite[Vgl.][S.~3--4]{butarbutar2023}

\subsection{Zielsetzung und Forschungsfragen}

Ziel der vorliegenden Arbeit ist es, auf Basis einer systematischen Literaturanalyse diejenigen Change-Management-Faktoren zu identifizieren und zu systematisieren, die den Erfolg von ERP-Implementierungen in mittelst\"andischen Unternehmen beeinflussen. Die Arbeit verbindet dazu zwei etablierte, bislang jedoch selten integrativ betrachtete Forschungsfelder: die ERP-Erfolgsforschung und das Change Management. W\"ahrend beide Felder f\"ur sich genommen umfangreich untersucht sind, fehlt es an Modellen, die beide Perspektiven systematisch verkn\"upfen. Insbesondere die Frage, welche Change-Management-Faktoren unter den besonderen Rahmenbedingungen mittelst\"andischer Unternehmen -- flache Hierarchien, knappe personelle und finanzielle Ressourcen sowie eine im Vergleich zu Gro\ss{}unternehmen geringere IT-Affinit\"at -- tats\"achlich erfolgskritisch sind, ist bislang nicht hinreichend beantwortet.

Daraus leitet sich die zentrale Forschungsfrage der Arbeit ab:

\begin{quote}
Welche Change-Management-Faktoren beeinflussen den Erfolg von ERP-Implementierungen in mittelst\"andischen Unternehmen?
\end{quote}

Zur strukturierten Beantwortung wird die Hauptforschungsfrage in vier Sub-Fragen zerlegt, die jeweils einen Teilaspekt des Untersuchungsgegenstands adressieren und sich unmittelbar in den Aufbau der Arbeit \"ubersetzen lassen.

\textbf{Sub-Frage 1 (SF1):} Welche Dimensionen des IS-Erfolgsmodells nach DeLone und McLean sind f\"ur den Kontext mittelst\"andischer Unternehmen besonders relevant und wie l\"asst sich der Erfolg einer ERP-Implementierung anhand dieses Modells operationalisieren?

Das Modell von DeLone und McLean\footcite[Vgl.][S.~9--10]{delone2003} gilt als eines der am besten validierten Erfolgsma\ss{}e der Informationssystemforschung, wurde jedoch \"uberwiegend in Gro\ss{}unternehmen und internationalen Kontexten gepr\"uft. Seine Anwendung auf ERP-Systeme im Mittelstand erfordert daher eine kritische Pr\"ufung der Dimensionen Systemqualit\"at, Informationsqualit\"at, Servicequalit\"at, Nutzerzufriedenheit und Nettonutzen sowie ihrer relativen Gewichtung. SF1 wird in Kapitel~2.2 theoretisch fundiert und liefert die Erfolgsdimensionen, auf die die Analyse in Kapitel~4.3 zur\"uckgreift.

\textbf{Sub-Frage 2 (SF2):} Welche Change-Management-bezogenen Erfolgsfaktoren f\"ur ERP-Implementierungen lassen sich aus der Forschungsliteratur ab 2015 identifizieren?

Der zeitliche Rahmen ab 2015 ist bewusst gew\"ahlt: Durch Digitalisierung, Cloud-L\"osungen und agile Vorgehensmodelle haben sich die Anforderungen an ERP-Projekte grundlegend ver\"andert, sodass \"altere, von monolithischen Einf\"uhrungsprojekten gepr\"agte Studien nur noch bedingt auf die heutige Praxis \"ubertragbar sind. Die Faktoren werden im Rahmen der systematischen Literaturanalyse extrahiert, kategorisiert und hinsichtlich ihrer empirischen Fundierung bewertet. Die Beantwortung von SF2 erfolgt in Kapitel~4.2.

\textbf{Sub-Frage 3 (SF3):} Welche Rolle spielen die in der Literatur als besonders erfolgskritisch geltenden Faktoren -- insbesondere Kommunikation, F\"uhrungsunterst\"utzung und Schulungsma\ss{}nahmen -- im Rahmen von ERP-Implementierungen in mittelst\"andischen Unternehmen und wie lassen sich diese anhand des ADKAR-Modells bewerten?

F\"ur diese Faktoren fehlen bislang detaillierte Untersuchungen dazu, wie sie speziell im Mittelstand wirken und welche Auspr\"agungen sie annehmen m\"ussen, um nachhaltig zum Projekterfolg beizutragen. Das ADKAR-Modell\footcite[Vgl.][]{hiatt2006} bietet hierf\"ur einen strukturierten Bewertungsrahmen, der die f\"unf Phasen des individuellen Ver\"anderungsprozesses abbildet und so eine differenzierte Analyse erm\"oglicht. SF3 wird auf Basis der Befunde aus Kapitel~4.2 in Verbindung mit dem in Kapitel~2.3 eingef\"uhrten ADKAR-Modell beantwortet und in Kapitel~4.3 mit den Erfolgsdimensionen verkn\"upft.

\textbf{Sub-Frage 4 (SF4):} Inwiefern lassen sich aus den Befunden der Literaturanalyse konkrete Handlungsempfehlungen f\"ur die Planung und Durchf\"uhrung von Change-Management-Ma\ss{}nahmen in ERP-Einf\"uhrungsprojekten mittelst\"andischer Unternehmen ableiten?

Die Arbeit verfolgt neben dem wissenschaftlichen Erkenntnisgewinn ein anwendungsorientiertes Ziel: Die abgeleiteten Empfehlungen sollen Entscheidungstr\"agern in mittelst\"andischen Unternehmen eine Orientierung bieten und dazu beitragen, die hohen Misserfolgsquoten von ERP-Projekten zu reduzieren. SF4 wird in Kapitel~5 beantwortet.

Die Zusammenf\"uhrung der Teilergebnisse und die abschlie\ss{}ende Beantwortung der Hauptforschungsfrage erfolgen in Kapitel~6.1.

Einen \"Uberblick \"uber die Zuordnung der Forschungsfragen zu den Kapiteln gibt Tabelle~\ref{tab:ff-mapping}.

\begin{table}[H]
\centering
\begin{tabularx}{\textwidth}{lX}
\toprule
\textbf{Forschungsfrage} & \textbf{Beantwortung in Kapitel} \\
\midrule
SF1: D\&M-Dimensionen im Mittelstand & 2.2 (Theorie) / 4.3 (Analyse) \\
SF2: Identifikation der CM-Faktoren & 4.2 (Ergebnisse) \\
SF3: Kommunikation, F\"uhrung, Schulung im ADKAR-Modell & 2.3 (Theorie) / 4.2--4.3 (Analyse) \\
SF4: Ableitung von Handlungsempfehlungen & 5 (Empfehlungen) \\
Hauptfrage: Synthese der Ergebnisse & 6.1 (Fazit) \\
\bottomrule
\end{tabularx}
\caption{Zuordnung der Forschungsfragen zu den Kapiteln}
\label{tab:ff-mapping}
\end{table}

\subsection{Aufbau der Arbeit}

Die Arbeit beginnt mit den theoretischen Grundlagen (Kapitel~2), in denen die relevanten Fachkonzepte -- ERP-Systeme, das IS-Erfolgsmodell nach DeLone \& McLean sowie das ADKAR-Modell -- aufbereitet und in einem vorl\"aufigen Bezugsrahmen zusammengef\"uhrt werden. Kapitel~3 stellt das methodische Vorgehen dar, das die systematische Literaturanalyse, die Suchstrategie und den Auswahlprozess nach PRISMA transparent macht. Kapitel~4 pr\"asentiert die Ergebnisse der systematischen Literaturanalyse, die in Kapitel~5 zu Handlungsempfehlungen f\"ur den Mittelstand verdichtet werden. Die Arbeit schlie\ss{}t mit Fazit und Ausblick (Kapitel~6).
